\section*{अध्याय \ref{ch:b3:identical-particles} --- सर्वसम कण}

\begin{solution}{exo:b3:identical-particles:1}
(क) दो बार अदला-बदली कुछ न करने के बराबर है: \omterm{def:b3:quantum-formalism:observable}{अभिलक्षणिक मान} $\pm1$। (ख)
\omterm{def:b3:identical-particles:exchange}{सर्वसम कण} $\hat H$ में सममित रूप से आते हैं, अतः $\hat H\hat
P_{12} = \hat P_{12}\hat H$: सममिति-वर्ग कभी नहीं बदलता। (ग)
$\ket a\otimes\ket b$ और $\ket b\otimes\ket a$ एक ही भौतिक स्थिति का
वर्णन करने वाली \emph{भिन्न} अवस्थाएँ होतीं --- एक अप्रेक्ष्य भेद, जिसे
संरूपण को ढोना ही नहीं चाहिए। (घ)
$^1$H (2 फ़र्मिऑन): बोसॉन; ड्यूटीरियम (3): फ़र्मिऑन; $^4$He: बोसॉन;
$^3$He: फ़र्मिऑन; H$_2$ (4): बोसॉन; फ़ोटॉन: बोसॉन।
\end{solution}

\begin{solution}{exo:b3:identical-particles:2}
(क) $(\ket a\otimes\ket b - \ket b\otimes\ket a)/\sqrt2$। (ख)
सर्वसमतः शून्य। (ग) छह क्रमचयों पर बदलते चिह्नों वाला योग --- कणों के
सामने कक्षकों का $3\times3$ सारणिक; दो पंक्तियों की अदला-बदली सारणिक
का चिह्न पलट देती है। (घ) जिस सारणिक के दो स्तंभ बराबर हों वह लुप्त हो
जाता है: एक ही कक्षक में दो फ़र्मिऑन एक दोहराया हुआ स्तंभ हैं।
\end{solution}

\begin{solution}{exo:b3:identical-particles:3}
(क) $3^2 = 9$। (ख) $3$ युगल $+$ $3$ जोड़े $= 6$। (ग) $3$ जोड़े।
(घ) छह एक-कण अवस्थाएँ (तीन कक्षक $\times$ दो चक्रण):
$\binom62 = 15$।
\end{solution}

\begin{solution}{exo:b3:identical-particles:4}
(क) C: $1s^22s^22p^2$; N: $2p^3$; O: $2p^4$; Ne: $2p^6$। (ख)
नाइट्रोजन के तीन $2p$ इलेक्ट्रॉन तीनों कक्षकों को एक-एक करके घेरते
हैं, चक्रण समांतर। (ग) समांतर चक्रण प्रतिसममित स्थानिक अवस्था थोपते
हैं, जिसका विनिमय विवर कूलॉम प्रतिकर्षण बचाता है --- संरेखण
स्थिरवैद्युत बचत से खरीदा जाता है। (घ) हीलियम (अपने उत्कृष्ट कुल के
साथ); और फ़्लुओरीन, जो नियॉन की बंदी से एक इलेक्ट्रॉन कम है, वह
इलेक्ट्रॉन लगभग किसी से भी छीन लेता है।
\end{solution}

\begin{solution}{exo:b3:identical-particles:5}
(क) बोसॉन: $2E_1$; चक्रणहीन फ़र्मिऑन: $E_1 + 4E_1 = 5E_1$;
इलेक्ट्रॉन: $2E_1$ (विपरीत चक्रण $n = 1$ साझा करते हैं)। (ख) बोसॉन और
इलेक्ट्रॉन: $E_1 + 4E_1 = 5E_1$; चक्रणहीन फ़र्मिऑन: $E_1 + 9E_1 =
10E_1$। (ग) फ़र्मिऑनों को दबाने से ऊँचे-से-ऊँचे स्तर भरने पड़ते हैं:
प्रति दबाव ऊर्जा-लागत एक बहिर्मुखी दाब है, जिसे कोई शीतलन नहीं हटाता
--- भ्रूण रूप में पदार्थ की असंपीड्यता। (घ)
$\sum_{n=1}^N n^2 \approx N^3/3$: फ़र्मी सागर का कुल
$N^3$ की तरह बढ़ता है।
\end{solution}

\begin{solution}{exo:b3:identical-particles:6}
(क) दोनों पद $\vect r_1 = \vect r_2$ पर कट जाते हैं। (ख) प्रसार करने
पर आड़े पद $\mp2|\bra a\hat{\vect r}\ket b|^2$ का योगदान देते हैं:
प्रतिसममित अवस्था का माध्य वर्ग पृथक्करण \emph{बड़ा} है। (ग) अधिक दूर
होने का अर्थ है कम प्रतिकर्षण ऊर्जा: $E_- <
E_+$। (घ) हीलियम का त्रिक (प्रतिसममित समष्टि) उसी विन्यास के एकल से
नीचे पड़ता है --- ऑर्थोहीलियम पैराहीलियम के नीचे, बची हुई
\qty{0.8}{eV} प्रतिकर्षण जितना।
\end{solution}

\begin{solution}{exo:b3:identical-particles:7}
(क) $1s^2$ एक सममित स्थानिक अवस्था है: अतः चक्रणों को प्रतिसममित एकल
बनाना ही होगा। (ख) विपाटन $= 2J$: $J \approx
\qty{0.4}{eV}$। $a_0$ पर सच्ची चुंबकीय द्विध्रुव--द्विध्रुव ऊर्जा:
$\mu_0\mu_{\text{B}}^2/4\pi a_0^3 \approx \qty{4e-4}{eV}$
--- हज़ार गुना छोटी: वह ``चक्रण बल'' दरअसल चक्रण का रंग ओढ़े हुए
स्थिरवैद्युतिकी है। (ग) $1\,^1S$ तक क्षय के लिए एक चक्रण-पलटाव
(एकल--त्रिक) \emph{और} एक $s \to s$ संक्रमण ($\Delta l =
0$) चाहिए: दोहरा वर्जित। (घ) $8000$ s बनाम \qty{1.6}{ns}: दो
चयन-नियमों से बारह कोटियाँ --- अर्धस्थायित्व नियमों के ढेर से बनता है,
बलों के कमज़ोर पड़ने से नहीं।
\end{solution}

\begin{solution}{exo:b3:identical-particles:8}
(क) कुल विनिमय चिह्न $=$ (चक्रण भाग) $\times (-1)^J$ होना चाहिए
$-1$: त्रिक (सममित चक्रण) विषम $J$ लेता है, और एकल सम $J$।
(ख) बहुलताओं से $3:1$। (ग) 3/4 ऑर्थो-अंश
पैरा ($J = 0$) तक शिथिल होता है, और प्रति अणु $2B \approx \qty{15}{meV}$
छोड़ता है --- जो द्रव की गुप्त ऊष्मा से \emph{अधिक} है:
अनुत्प्रेरित टंकी धीरे-धीरे उबलकर सूख जाती है। संयंत्र द्रवण के दौरान
उत्प्रेरक पर ऑर्थो $\to$ पैरा बदल देते हैं। (घ) द्रव
हाइड्रोजन की रॉकेटरी --- हर प्रक्षेपण-मंच को नाभिकीय-चक्रण
सांख्यिकी की एक समस्या विरासत में मिलती है।
\end{solution}

\begin{solution}{exo:b3:identical-particles:9}
(क) दरें आयाम के वर्ग की तरह चलती हैं: किसी अधिभुक्त मोड में $n + 1$,
और उलटी दिशा में $n$। (ख) $10^{10} + 1$ बनाम $1$: भीड़ भरा मोड दस
अरब गुने से जीतता है। (ग) पसंदीदा मोड में एक स्वतःस्फूर्त फ़ोटॉन
पासा झुका देता है; और हर उद्दीपित फ़ोटॉन
$n$ बढ़ाकर झुकाव तीखा करता है: सर्वसम फ़ोटॉनों का हिमस्खलन ---
एक ही मोड में सिमटती वृद्धि। (घ) $\hat a^\dagger\ket n =
\sqrt{n+1}\,\ket{n+1}$ में: बोसॉनी मोड ही दोलक है, और
वह संवर्धन उसका सोपान गुणांक है।
\end{solution}

\begin{solution}{exo:b3:identical-particles:10}
(क) हर गैस दुगने आयतन में फैलती है: प्रति गैस $\Delta S =
Nk_{\text{B}}\ln2$। (ख) टोंटी बंद करके फिर खोलना बिना किसी
अवस्था-परिवर्तन के एन्ट्रॉपी गढ़ देता ---
हिसाब की सतत ठगी। (ग)
नाम वाले अणुओं के साथ मिली हुई गैस के पास $\binom{2N}{N}$ अतिरिक्त
``किस ओर'' नियतन होते हैं, जिनका मोल $2Nk_{\text{B}}\ln2$ है; और सारी
अवस्था-गणनाओं को $N!$ से भाग देना ---
सर्वसम कणों के लिए ईमानदार गणना --- ठीक इसी को हटा देता है।
(घ) वही $N!$, जिसकी चिरसम्मत भौतिकी को ज़रूरत थी और जिसे वह उचित नहीं
ठहरा सकती थी, दूर से देखा गया सममितीकरण अभिगृहीत है।
\end{solution}

\begin{solution}{exo:b3:identical-particles:11}
(क) $E_{\text{F}} = \hbar^2(3\pi^2n_{\text{e}})^{2/3}/2m_{\text{e}}
\approx \qty{0.16}{MeV}$ --- इलेक्ट्रॉन की \omterm{def:b3:relativistic-dynamics:momentum-energy}{विराम
ऊर्जा} का एक तिहाई: तारा आपेक्षिकीय प्रांत के किनारे बैठा है। (ख)
$k_{\text{B}}T \approx \qty{0.9}{keV} \ll E_{\text{F}}$: फ़र्मी
पैमाने पर वह श्वेत-तप्त तारा गहरे अपह्रासित है --- अर्थात् ``ठंडा''।
(ग) इलेक्ट्रॉन पाउली-शैली में
$\qty{0.16}{MeV}$ तक चुनते जाते हैं; तारे को दबाने से हर अधिभुक्त
स्तर ऊपर उठता है, और वह ऊर्जा-प्रवणता ताप से स्वतंत्र एक दाब है ---
अपवर्जन, ऊष्मा नहीं, बोझ ढोता है। (घ)
आपेक्षिकीय अपह्रासिता दाब-नियम को तब तक नरम करती है जब तक किसी परिमित
द्रव्यमान पर गुरुत्व न जीत जाए: चंद्रशेखर सीमा, जो
\cref{ch:b3:astrophysics} के लिए रखी है।
\end{solution}

\begin{solution}{exo:b3:identical-particles:12}
(क) $\hat{\vect S}_1\cdot\hat{\vect S}_2 = \tfrac12(\hat S^2 -
\tfrac32\hbar^2)$: एकल पर $-\tfrac34\hbar^2$,
और त्रिक पर $+\tfrac14\hbar^2$ --- कथित हैमिल्टनियन
$2J$ का विपाटन फिर से दे देता है। (ख) संरेखित पड़ोसियों का कोई जालक:
स्वतःस्फूर्त चुंबकन --- अर्थात् लौहचुंबकत्व। (ग) $J/k_{\text{B}}
\approx \qty{1200}{K}$: लोहे का \qty{1043}{K} ठीक इसी
पैमाने का है। (घ) द्विध्रुव--द्विध्रुव चुंबकत्व $10^{-4}$ eV का है:
वह मिलिकेल्विन पर क्रमित होता। लौहचुंबकत्व कूलॉम प्रतिकर्षण है, जो
विनिमय-ज्यामिति के ज़रिए चक्रण-विन्यास चुनता है।
\end{solution}

\begin{solution}{pb:b3:identical-particles:1}
\textbf{1.} कक्षक $(n, l, m)$, जिनमें $n^2$ प्रति $n$; चक्रण दुगुना
करता है; और पाउली प्रति पूर्ण नामपत्र एक इलेक्ट्रॉन की सीमा बाँधता
है: $2$, $8$, $18$।
\textbf{2.} भीतरी इलेक्ट्रॉन नाभिक को भिन्न $l$ के लिए भिन्न ढंग से
परिरक्षित करते हैं: भेदक $s$ कक्षक कम भेदक $p$ और $d$ से नीचे बैठ
जाते हैं --- ऊर्जा-क्रम $(n, l)$ का अनुसरण करता है, अकेले $n$ का
नहीं।
\textbf{3.} क्रम $1s\,2s\,2p\,3s\,3p\,4s\,3d$; Na: [Ne]$3s^1$;
Cl: [Ne]$3s^23p^5$; K: [Ar]$4s^1$; Fe: [Ar]$4s^23d^6$।
\textbf{4.} पंक्तियाँ हर उत्कृष्ट विन्यास पर बंद होती हैं: $2$
($1s$), $8$ ($2s2p$), फिर $8$ ($3s3p$ --- क्योंकि $3d$ $4s$ से ऊपर
पड़ता है और प्रतीक्षा करता है), फिर $18$ ($4s3d4p$)।
\textbf{5.} एक बंद कोश: अगले कक्षक तक बड़ा अंतराल, देने के लिए कोई
सस्ता इलेक्ट्रॉन नहीं, लेने के लिए कोई जगह नहीं।
\textbf{6.} $Z_{\text{प्र}}$ नाभिकीय आवेश में से भीतरी (और आंशिक रूप
से, उसी कोश के) इलेक्ट्रॉनों का परिरक्षण घटाकर मिलता है: सोडियम का
संयोजी इलेक्ट्रॉन $11 - 10 \approx 1$ देखता है; और हीलियम के दोनों
नंगे $Z = 2$ को बाँटते हैं, प्रत्येक दूसरे को केवल एक तिहाई परिरक्षित
करता है।
\textbf{7.} Li $\to$ Be तक चढ़ती है, B पर गिरती है; N $\to$ O पर
गिरती है: यही दोनों झुकाव।
\textbf{8.} बोरॉन $2p$ खोलता है, जो $2s$ से ऊँचा और कम भेदक है: उसका
नया इलेक्ट्रॉन बस हटाने में आसान है।
\textbf{9.} ऑक्सीजन को पहली बार एक ही $p$ कक्षक में दो इलेक्ट्रॉन
\emph{युग्मित} करने पड़ते हैं: अतिरिक्त प्रतिकर्षण, खोया हुआ
विनिमय-स्थायित्व --- चौथे $2p$ इलेक्ट्रॉन पर छूट लग जाती है।
\textbf{10.} हर इलेक्ट्रॉन दूसरे को अपूर्ण रूप से परिरक्षित करता है:
अनुरूपण $24.6 = 13.6\,Z_{\text{प्र}}^2$ देता है $Z_{\text{प्र}} =
1.34$ --- अर्थात् दो-तिहाई आवेश परिरक्षित।
\textbf{11.} लगभग पूरा क्रोड-परिरक्षण ($Z_{\text{प्र}}
\approx 1.3$) और बाहरी इलेक्ट्रॉन का $n = 2$: दोनों बंधन-ऊर्जा घटा
देते हैं।
\textbf{12.} जुड़ते इलेक्ट्रॉन \emph{उसी} कोश में जाते हैं और
एक-दूसरे को खराब परिरक्षित करते हैं, जबकि $Z$ चढ़ता है:
$Z_{\text{प्र}}$ बढ़ता है और पूरा कोश सिकुड़ जाता है --- नियॉन
लीथियम से छोटा है।
\textbf{13.} किसी उत्कृष्ट क्रोड के ऊपर एक सस्ता $s$ इलेक्ट्रॉन:
आयनन $\sim\qty{5}{eV}$, तत्काल $+1$ आयन, प्रचंड इलेक्ट्रॉन-दान
--- कुल-लक्षण एक ही विन्यास हैं।
\textbf{14.} किसी $p^6$ बंदी में एक विवर: विशाल इलेक्ट्रॉन
बंधुता, $-1$ आयन, और विवरों का युग्मन X$_2$ अणुओं में।
\textbf{15.} $2s$ और $2p$ पास-पास पड़ते हैं: $2s^22p^2 \to$ चार
अर्ध-भरे संकरों तक पदोन्नति कम लागत की है और दो अतिरिक्त बंध खरीद
देती है --- कार्बन की चतुःसंयोजकता एक भुनाई गई निकट-अपह्रासिता है।
\textbf{16.} चार इलेक्ट्रॉनों के साथ बंधक \emph{और} प्रतिबंधक दोनों
संयोजन भर जाते हैं: लाभ कट जाते हैं, कोई सहसंयोजी बंध नहीं ---
हीलियम अकेला ही रहता है।
\textbf{17.} $d$ इलेक्ट्रॉन स्वयं को $4s$ की खाल के भीतर गाड़ लेते
हैं: बाहरी रसायन, जिसे $s$ कोश तय करता है, दस तत्त्वों में बहुत कम
बदलता है --- एक रासायनिक पठार।
\textbf{18.} स्कैंडियम, [Ar]$4s^23d^1$: पहला अयुग्मित $d$
इलेक्ट्रॉन।
\textbf{19.} तीन संरेखित चक्रणों का एक चुंबकीय आघूर्ण
($\sim3\mu_{\text{B}}$); और उस संरेखण की कीमत विनिमय चुकाता है ---
कूलॉम बचत, चुंबकीय बल नहीं।
\textbf{20.} थोड़े-से विभेद्य इलेक्ट्रॉन सब के सब $1s$ की ओर भीड़
लगा सकते (सममितीकृत अवस्थाओं का ठीक-ठीक कटाव न होने से): कोश रिसने
लगते, बंदियाँ धुँधली पड़तीं, और आवर्त नियम घुल जाता।
\textbf{21.} $v/c \approx Z\alpha = 0.58$: आपेक्षिकीय
द्रव्यमान-वृद्धि $6s$ कक्षक को सिकोड़कर स्थायी बनाती है; सोने का
$5d \to 6s$ अवशोषण नीले में सरक जाता है (और परावर्तित पीला छोड़ देता
है), तथा पारे का कसा हुआ $6s^2$ युग्म इतना कमज़ोर बंध बनाता है कि वह
धातु द्रव है।
\textbf{22.} वे लगभग बराबर चलती हैं: अधिभोग पर निर्भर क्रम,
$4s$ का पहले आयनित होना, और Cr तथा Cu की $3d^5$/$3d^{10}$ उधारियाँ
--- अर्ध-भरे और भरे उपकोश विनिमय-स्थायित्व सस्ते में खरीद लेते हैं।
\textbf{23.} रसायन इलेक्ट्रॉनों का मामला है, और इलेक्ट्रॉन केवल
नाभिकीय \emph{आवेश} अनुभव करते हैं: द्रव्यमान बस नन्हे कंपन
(समस्थानिक) खिसकावों से भीतर आता है।
\textbf{24.} कोश-संरचना (जो उत्कृष्ट-गैस तर्क का पक्ष लेती है) बनाम
आपेक्षिकीय स्थायीकरण और $Z\alpha \to 1$ पर निरा कूलॉम तनाव:
अतिभारी रसायन उन्हीं की रस्साकशी है।
\textbf{25.} कक्षकों की एक सूची, एक दुगुनापन, एक ऋण चिह्न: पंक्तियाँ
$2, 8, 8, 18$, बोरॉन और ऑक्सीजन पर झुकाव, क्षार से उत्कृष्ट तक के
कुल, \qty{5.1}{eV} से \qty{24.6}{eV} तक की आयनन-सीढ़ी --- आवर्त नियम
व्युत्पन्न हुआ, रटा नहीं गया।
\end{solution}
